% ============================================================================
% Doc. 264 -- Ferrotoroidicity and Quaternary Memory: All Formulas
%             Reinterpreted from the FFGFT Perspective
%
% Reference paper: Qureshi et al., Nature Communications 17:4033 (2026)
% DOI: 10.1038/s41467-026-70767-8
%
% Status: Reduction layer -- these are structural analogies and bridge
% arguments, not core derivations. Where FFGFT provides a natural home
% for a formula, this is noted; where the connection is speculative,
% it is marked explicitly.
% ============================================================================

\section*{Preamble: Why This Paper Matters for FFGFT}

Qureshi et al.\ (2026, \emph{Nature Communications} 17:4033) report the first
experimental demonstration of quaternary (4-state) non-volatile memory in a
purely antiferromagnetic single crystal of LiNi$_{0.8}$Fe$_{0.2}$PO$_4$, using
\emph{ferrotoroidic} domain control via perpendicular electric and magnetic
fields. The four distinct magnetic domains arise from a spontaneous spin
rotation below a phase transition temperature $T_4 = 21$\,K, and their
selective population is verified by Spherical Neutron Polarimetry (SNP).

From the FFGFT perspective, this paper is remarkable for three independent
reasons:
\begin{enumerate}
  \item The \emph{torus} is not only the fundamental geometric object of
        FFGFT (T$^4$, Dok.~001/002) but appears directly as the physical
        origin of the 4-state structure in the experiment: four domain states
        = two toroidal directions $\times$ two orientation states, a $2\times2$
        decomposition that mirrors the mode labelling on a compact torus.
  \item The \emph{polarization matrix formalism} (Blume-Maleev equations) is
        a rotation-matrix representation of state evolution -- structurally
        identical to the unitary operators on the T$^4$ mode space in FFGFT
        (Dok.~203, Recursion $\mathcal{R}$).
  \item The \emph{ferrotoroidic order parameter} breaks both space-inversion
        and time-inversion symmetry simultaneously -- exactly the symmetry
        structure of the FFGFT winding modes, which carry both a spatial
        winding number and a temporal orientation.
\end{enumerate}

\noindent\textbf{Status of this document.} Every formula from the paper is
reproduced and reinterpreted. The FFGFT interpretations are reduction layer
unless explicitly marked as bridge or core. The paper's own results are not
contested; the reinterpretation adds a geometric reading, not a physical
correction.

% ============================================================================
\section{Toroidization Formula (Eq.~1 in the Paper)}
% ============================================================================

\subsection*{The formula}

\begin{equation}
  \mathbf{T} = 2S\,\frac{\varepsilon_a}{\Omega}\cos\varphi \cdot \mathbf{b}
  \times \mathbf{c}
  \tag{Paper~Eq.~1}
\end{equation}

\noindent where $S$ is the spin magnitude, $\varepsilon_a/\Omega$ is a
geometric prefactor (unit cell volume $\Omega$, spin displacement $\varepsilon_a$),
$\varphi$ is the rotation angle of the magnetic moments in the $a$--$b$ plane,
and $\mathbf{b}\times\mathbf{c}$ defines the toroidal direction along the
crystallographic $c$ axis.

\subsection*{FFGFT interpretation}

In FFGFT, the torus T$^4$ has four compact dimensions. The toroidization
$\mathbf{T}$ is the macroscopic projection of a \emph{winding flux}:
\begin{itemize}
  \item $S$ corresponds to the winding quantum number $n$ of a torus mode
        (Dok.~001, fundamental relation T$\tilde{\cdot}m=1$): the spin
        magnitude measures how many times the mode winds around the compact
        direction.
  \item $\varepsilon_a/\Omega$ is the \emph{packing density} of the winding
        in the unit cell -- directly analogous to the packing deficit
        $\xi = 4/30000$ in FFGFT, which measures how much the 3D packing
        of space deviates from a smooth manifold (Dok.~009).
  \item $\cos\varphi$ is the \emph{projection} of the winding onto the
        observable direction. In FFGFT, this corresponds to the
        decompactification projection: the full T$^4$ winding is
        unit-free (Layer~1), and $\cos\varphi$ represents the fraction
        that becomes observable in the SI-projected (Layer~2) description.
  \item $\mathbf{b}\times\mathbf{c}$ defines the toroidal axis -- in FFGFT,
        this is the projection direction from the 4D compact space onto the
        3D observable subspace. The cross product structure is the same as
        in the T$^4$ volume form restricted to a 2-cycle.
\end{itemize}

\noindent\textbf{Key observation:} $\mathbf{T}$ is independent of the
rotation angle $\varphi$ in the sense that its \emph{direction} (sign) does
not change as $\varphi$ varies -- only its magnitude changes. This is
structurally identical to the FFGFT non-closure theorem (Dok.~193): the
topological charge (winding number) is invariant under continuous
deformations of the mode within a winding class.

% ============================================================================
\section{Polarization Matrix (Eq.~2 and Eq.~3)}
% ============================================================================

\subsection*{The formula}

\begin{equation}
  \mathbf{P}_f = \mathbf{P}\,\mathbf{P}_i
  \tag{Paper~Eq.~2}
\end{equation}

\begin{equation}
  \mathcal{P} = \frac{1}{I}
  \begin{pmatrix}
    N^2 - M_\perp^2 & J_{nz} & J_{ny} \\
    J_{nz} & N^2 + M_{\perp y}^2 - M_{\perp z}^2 & R_{yz} \\
    J_{ny} & R_{yz} & N^2 - M_{\perp y}^2 + M_{\perp z}^2
  \end{pmatrix}
  \tag{Paper~Eq.~3}
\end{equation}

\noindent where $N$ is the nuclear structure factor, $M_\perp$ is the magnetic
interaction vector (perpendicular component of the magnetic structure factor),
$J_{ni} = 2\,\mathrm{Im}(NM_{\perp i}^*)$ are nuclear-magnetic interference
terms, and $R_{yz} = 2\,\mathrm{Re}(M_{\perp y}M_{\perp z}^*)$.

\subsection*{FFGFT interpretation}

The transformation $\mathbf{P}_f = \mathbf{P}\,\mathbf{P}_i$ is a
\emph{rotation matrix acting on the neutron spin state}. In FFGFT, this is
the direct analogue of the recursion operator $\mathcal{R}$ (Dok.~203) acting
on the T$^4$ mode vector: the state vector transforms under a unitary rotation
determined by the geometric structure of the scattering object.

The matrix $\mathcal{P}$ has four distinct forms for the four domains -- each
domain corresponds to a \emph{different winding configuration} on the torus,
and each winding configuration produces a different rotation of the mode
vector. This is structurally identical to the four winding classes in FFGFT:
two toroidal directions (time-reversal pairs, $\pm$ winding) $\times$ two
orientation states (spatial reflection pairs), giving a $\mathbb{Z}_2 \times
\mathbb{Z}_2$ group structure.

The nuclear-magnetic interference terms $J_{ni}$ are non-zero because the
magnetic structure factor $M_\perp$ is \emph{purely imaginary} in the
antiferromagnetic phase (the moments related by inversion are antiparallel).
In FFGFT, imaginary structure factors correspond to modes that are
\emph{anti-phase} on opposite sides of the torus -- precisely the
antiferromagnetic (head-to-tail) configuration that defines the ferrotoroidic
ground state.

The diagonal element $N^2 - M_\perp^2$ measures the imbalance between nuclear
(scalar) and magnetic (vector) scattering -- in FFGFT, this is the imbalance
between the scalar field $E$ (Lagrangian $\mathcal{L}=\xi(\partial E)^2$,
Dok.~010) and the vector field components, which determines the energy density
$T_{00} = \xi[(\partial_t E)^2 + (\nabla E)^2]$.

% ============================================================================
\section{General Polarization Matrix Including Beam Imperfections (Eq.~9)}
% ============================================================================

\subsection*{The formula}

\begin{equation}
  \mathcal{P} =
  \begin{pmatrix}
    \frac{p_{fx}}{p_{ix}}\frac{N^2-M_\perp^2-J_{yz}}{I_x} &
    \frac{p_{fx}}{p_{iy}}\frac{J_{nz}-J_{yz}}{I_y} &
    \frac{p_{fx}}{p_{iz}}\frac{J_{ny}-J_{yz}}{I_z} \\[6pt]
    \frac{p_{fy}}{p_{ix}}\frac{J_{nz}+R_{ny}}{I_x} &
    \frac{p_{fy}}{p_{iy}}\frac{N^2+M_{\perp y}^2-M_{\perp z}^2+R_{ny}}{I_y} &
    \frac{p_{fy}}{p_{iz}}\frac{R_{yz}+R_{ny}}{I_z} \\[6pt]
    \frac{p_{fz}}{p_{ix}}\frac{J_{ny}+R_{nz}}{I_x} &
    \frac{p_{fz}}{p_{iy}}\frac{R_{yz}+R_{nz}}{I_y} &
    \frac{p_{fz}}{p_{iz}}\frac{N^2-M_{\perp y}^2+M_{\perp z}^2+R_{nz}}{I_z}
  \end{pmatrix}
  \tag{Paper~Eq.~9}
\end{equation}

\noindent where $p_{fx}, p_{fy}, p_{fz}$ are the final analysis efficiencies
and $p_{ix}, p_{iy}, p_{iz}$ are the initial polarization values along the
three reference axes; $R_{ni} = 2\,\mathrm{Re}(NM_{\perp i}^*)$,
$J_{yz} = 2\,\mathrm{Im}(M_{\perp y}M_{\perp z}^*)$, and the scattered
intensities $I_x, I_y, I_z$ now include polarization-dependent corrections.

\subsection*{FFGFT interpretation}

The factors $p_{fx}/p_{ix}$ etc.\ are \emph{projection efficiencies} -- they
measure how accurately the measurement apparatus resolves the polarization
state. In FFGFT, these correspond to the \emph{decompactification factors}:
the ratio between the observable (SI-projected, Layer~2) quantity and the
underlying geometric (unit-free, Layer~1) quantity. Perfect measurement
($p_{fi} = 1$) corresponds to the Layer~1 limit (no embedding price); real
measurement ($p_{fi} < 1$) corresponds to the Layer~2 projection with its
inherent incompleteness.

The quantities $R_{ni} = 2\,\mathrm{Re}(NM_{\perp i}^*)$ are
\emph{real} nuclear-magnetic interference terms, while $J_{ni}$ are imaginary
ones. In FFGFT, this separation is exactly the decomposition into
\emph{even} (cosine, time-symmetric) and \emph{odd} (sine, time-asymmetric)
winding mode components (Dok.~203, Recursion $\mathcal{R}$). The real
part $R_{ni}$ involves modes that are symmetric under time reversal; the
imaginary part $J_{ni}$ involves modes that change sign under time reversal
-- exactly the ferrotoroidic modes that break time-inversion symmetry.

The general intensity $I_i = N^2 + M_\perp^2 + p_{ii} J_{yz}$ (or similar
with $R_{ni}$) shows that the observable intensity depends on the
\emph{initial polarization} -- in FFGFT, this corresponds to the fact that
the observable energy density $T_{00}$ depends on the temporal component
$(\partial_t E)^2$: a field with no temporal variation has $T_{00}=0$
(Dok.~010), just as a completely unpolarized neutron beam ($p_{ii}=0$) gives
zero interference terms.

% ============================================================================
\section{Dzyaloshinskii--Moriya Hamiltonian (Eqs.~4--6)}
% ============================================================================

\subsection*{The formulas}

\begin{align}
  \mathcal{H}_1^{\mathrm{DM}} &= \mathbf{D}_{12}\cdot(\mathbf{S}_1\times\mathbf{S}_2)
    + \mathbf{D}_{34}\cdot(\mathbf{S}_3\times\mathbf{S}_4)
  \tag{Paper~Eq.~4} \\[6pt]
  \mathcal{H}_{1c}^{\mathrm{DM}} &= D_{12}^c
    (S_1^a S_2^b - S_1^b S_2^a - S_3^a S_4^b + S_3^b S_4^a)
  \tag{Paper~Eq.~5} \\[6pt]
  \mathcal{H}_{1b}^{\mathrm{DM}} &= -D_{12}^b
    (S_1^a S_2^c - S_1^c S_2^a - S_3^a S_4^c + S_3^c S_4^a)
  \tag{Paper~Eq.~6}
\end{align}

\noindent with $\mathbf{D}_{34} = -\mathbf{D}_{12}$ (antisymmetry under inversion),
and the components $D_{12}^b, D_{12}^c$ being the projections of the
Dzyaloshinskii--Moriya (DM) vector onto the crystallographic $b$ and $c$ axes.

\subsection*{FFGFT interpretation}

The Dzyaloshinskii--Moriya interaction is an \emph{antisymmetric exchange}:
$\mathbf{D}\cdot(\mathbf{S}_1\times\mathbf{S}_2)$ is odd under exchange of
spins 1 and 2 (compare: the symmetric Heisenberg exchange
$J\,\mathbf{S}_1\cdot\mathbf{S}_2$ is even). In FFGFT, antisymmetric
interactions correspond to the \emph{non-commutative} products of winding
modes -- the cross-product structure $S_i\times S_j$ is precisely the
structure that the Non-Closure Theorem (Dok.~193) identifies as the
topologically constrained product of mode labels from different winding
classes.

The relation $\mathbf{D}_{34} = -\mathbf{D}_{12}$ (due to inversion symmetry)
corresponds to the FFGFT symmetry that modes related by spatial inversion
have opposite winding numbers: if mode $n$ winds clockwise, the inversion-
related mode $-n$ winds anti-clockwise. This is the T$^4$ analogue of the
crystallographic inversion operator.

The components $D^c$ and $D^b$ (with $|D^b| < |D^c|$ from the geometry,
as shown in Fig.~6 of the paper) reflect that the DM vector is not
aligned with any crystallographic axis -- it lies at $35°$ from the $c$
axis in the $b$--$c$ plane. In FFGFT, this corresponds to the
\emph{anisotropic ξ-coupling}: the geometric parameter $\xi$ enters
differently in different compact directions, with the dominant coupling
in the direction of the largest winding mode (here, the $c$-axis
ferrotoroidic direction).

The explicit form of $\mathcal{H}_{1c}^{\mathrm{DM}}$ (Eq.~5) is a
sum of four terms $\pm S_i^a S_j^b$ -- the $ab$-plane spin cross product.
In FFGFT language, this is the \emph{mixed-mode bilinear} that couples
the $a$-axis winding (orientation domain) to the $b$-axis winding
(primary ordered moment): exactly the coupling between two different T$^4$
compact directions that produces the $2\times2$ domain structure.

% ============================================================================
\section{Polarization Tensor Element Formula (Eq.~7)}
% ============================================================================

\subsection*{The formula}

\begin{equation}
  P_{fi} = \frac{n_{fi}^{++} - n_{fi}^{+-}}{n_{fi}^{++} + n_{fi}^{+-}}
  \tag{Paper~Eq.~7}
\end{equation}

\noindent where $n_{fi}^{++}$ ($n_{fi}^{+-}$) counts neutrons with initial
spin parallel to $i$ and final spin parallel (antiparallel) to $f$.

\subsection*{FFGFT interpretation}

This is the standard quantum-mechanical expectation value of the
polarization operator, normalized by the total count. In FFGFT, it
corresponds to the ratio of two mode populations: modes in the ``$+$''
winding state vs.\ modes in the ``$-$'' winding state, after scattering
(i.e., after the recursion operator $\mathcal{R}$ has acted once).

The normalization by $n^{++} + n^{+-}$ (total count) is the FFGFT
\emph{decompactification}: the absolute number of modes (a Layer~1
quantity) is not directly observable; only the ratio (a dimensionless,
Layer~1-compatible quantity) is. This is the same structure as the
FFGFT $T_{\mathrm{CMB}}$ relation: only the ratio $T_{\mathrm{CMB}}/E_\xi = (16/9)\xi^2$
is a pure geometric prediction; the absolute Kelvin value requires
the SI embedding.

% ============================================================================
\section{Final Polarization Including Created/Annihilated Components (Eq.~8)}
% ============================================================================

\subsection*{The formula}

\begin{equation}
  \mathbf{P}_f = \mathcal{P}'\,\mathbf{P}_i + \mathcal{P}''
  \tag{Paper~Eq.~8}
\end{equation}

\noindent where $\mathcal{P}'$ is the purely rotational part of the
polarization transformation and $\mathcal{P}''$ is the created or
annihilated polarization (non-zero when the scattering process itself
creates polarization).

\subsection*{FFGFT interpretation}

The decomposition into a \emph{rotational} part $\mathcal{P}'$ and a
\emph{creation} part $\mathcal{P}''$ maps directly onto the FFGFT
distinction between two types of mode transformation:

\begin{itemize}
  \item $\mathcal{P}'\,\mathbf{P}_i$: \emph{coherent} transformation --
        the input state is rotated without energy exchange. This corresponds
        to the FFGFT recursion $\mathcal{R}$ (Dok.~203) acting on an existing
        mode: the winding number is preserved, only the phase (orientation)
        changes.
  \item $\mathcal{P}''$: \emph{spontaneous} polarization -- new polarization
        is created by the scattering process itself. This corresponds to
        \emph{mode creation} in FFGFT: a new mode is generated at a T$^4$
        node (Dok.~204, Fractal Boundary Condition), with a winding number
        determined by the local geometry rather than the input state.
\end{itemize}

\noindent In the specific magnetic structure of LiNi$_{0.8}$Fe$_{0.2}$PO$_4$,
$\mathcal{P}''=0$ because the magnetic structure factors are purely imaginary
(as noted in the Methods section of the paper). In FFGFT, this corresponds to
the case where the T$^4$ geometry is \emph{static}: no new modes are created
spontaneously, and all transformations are purely rotational (unitary). This
is the T$^4$ static ground state -- precisely the FFGFT cosmological picture
(static universe, no spontaneous mode creation = no expansion).

% ============================================================================
\section{The Four-Domain Structure: $2\times2$ on T$^4$}
% ============================================================================

The paper's central experimental result is that four distinct magnetic
domains exist, controlled by two independent handles:
\begin{enumerate}
  \item \textbf{Toroidal domains} (2 states): selected by $\mathbf{E}\times\mathbf{H}$
        (parallel or antiparallel to $c$).
  \item \textbf{Orientation domains} (2 states): selected by $H_c$ (a small
        $c$-component of the magnetic field, via the DM energy).
\end{enumerate}

\noindent The four combined states ($a_4', a_4'', b_4', b_4''$) are related by:
\begin{itemize}
  \item Time-inversion $1'$: maps $a_4 \leftrightarrow b_4$ (toroidal flip).
  \item 2-fold screw axis $2_1\|y$: maps $'\leftrightarrow''$ (orientation flip).
\end{itemize}

\subsection*{FFGFT interpretation}

This $\mathbb{Z}_2\times\mathbb{Z}_2$ group structure is the direct
experimental realization of the T$^4$ winding mode classification in FFGFT:

\begin{itemize}
  \item The \emph{toroidal} $\mathbb{Z}_2$: corresponds to the two possible
        orientations of a closed winding loop on the T$^4$ (clockwise vs.\
        anti-clockwise = $\pm\mathbf{T}$). This is the fundamental ``arrow of
        time'' degree of freedom in FFGFT: the intrinsic time $\tilde{T}$
        can point in two directions on the compact torus (Dok.~230,
        Bell correlations as T$^4$ geometry).
  \item The \emph{orientation} $\mathbb{Z}_2$: corresponds to the two possible
        spatial orientations of the winding plane within the 4D torus. In
        FFGFT, this is the spatial parity of the mode -- whether the winding
        wraps in the $+c$ or $-c$ spatial direction.
\end{itemize}

\noindent The four combined states are thus:
\[
  \{+\tilde{T},+c\},\quad \{+\tilde{T},-c\},\quad \{-\tilde{T},+c\},\quad
  \{-\tilde{T},-c\}
\]
which is exactly the FFGFT mode labelling for a T$^4$ with one temporal and
one spatial compact direction contributing to the observable domain structure.

\noindent\textbf{WBE connection:} The four-domain structure has an
effective dimensionality $D_{\mathrm{eff}} = 2$ (two binary degrees of
freedom). The WBE exponent for this effective dimension is
$\alpha = D_{\mathrm{eff}}/(D_{\mathrm{eff}}+1) = 2/3$. This is neither
$w=-3/4$ (Onur's 3D network) nor $w=-1$ (ΛCDM), but an intermediate value
corresponding to the 2D nature of the toroidal domain control surface.
This suggests that different effective dimensions apply at different scales
and for different physical subsystems -- which is consistent with the
FFGFT picture that $\xi$ encodes a fractional departure from integer
dimension at all scales (Dok.~242, ``Die Skalenleiter'').

% ============================================================================
\section{Quaternary Memory and FFGFT Information Structure}
% ============================================================================

The paper shows that 8 quaternary units yield $4^8 = 65536$ combinations,
compared to $2^8 = 256$ for binary -- a factor of 256 increase. This is
the \emph{power-law increase of storage density} enabled by going from
binary to quaternary logic.

In FFGFT, this corresponds to the fundamental observation that the T$^4$
torus carries \emph{more information per mode} than a point or a circle:
each mode is labelled by four winding numbers (one per compact dimension),
not just one (binary). The storage capacity grows as $4^N$ (for N modes,
each with 4 states) vs.\ $2^N$ (binary), which is exactly the ratio between
the T$^4$ phase space and a product of circles.

The \emph{non-volatile} nature of the memory -- the domain populations are
stable in zero applied field, proven by measuring at $E=H=0$ -- corresponds
in FFGFT to the \emph{topological stability} of winding modes: a mode with
non-zero winding number cannot decay to zero winding without crossing an
energy barrier (the minimal energy of the T$^4$ boundary condition,
Dok.~204). This is why toroidal memory is non-volatile: the topological
charge is conserved unless a conjugate field large enough to overcome the
anisotropy barrier is applied.

% ============================================================================
\section{Relation to the FFGFT Parameter $\xi$}
% ============================================================================

The geometry of the paper involves three key length scales:
\begin{itemize}
  \item The unit cell dimension ($\sim 10$\,\AA): sets the scale of the
        individual winding loop.
  \item The domain size ($\sim$mm): the macroscopic coherence length of the
        toroidal order.
  \item The phase transition temperatures $T_4=21$\,K and $T_2=25$\,K:
        the energy scales for the two levels of domain structure.
\end{itemize}

\noindent In FFGFT, the ratio of the smallest and largest scales is governed
by $\xi$:
\[
  \frac{L_0}{L_{\mathrm{domain}}} \sim \xi^n
  \quad \text{for some integer or fractional } n.
\]
The exact value of $n$ is not determined within this document (reduction
layer), but the existence of a $\xi$-controlled hierarchy between the atomic
scale (unit cell) and the macroscopic scale (domain) is a structural
prediction of FFGFT.

\noindent\textbf{Energy ratio:} The ratio of the two phase transition
temperatures is $T_4/T_2 = 21/25 = 0.840$. In FFGFT, ratios of energy
scales are related to powers of $\xi$:
\[
  1 - \xi = 1 - \frac{4}{30000} \approx 0.99987
  \quad \text{(first order)}
\]
The observed ratio $0.840$ is closer to $1-\xi^{1/4} = 1 - 0.107 \approx
0.893$ (a rough estimate). A precise geometric derivation of this ratio
from the FFGFT T$^4$ geometry is an open question (reduction layer).

% ============================================================================
\section{Summary: What the Experiment Reveals about the T$^4$ Structure}
% ============================================================================

\begin{itemize}
  \item \textbf{Toroidal order is real and controllable:} The experiment
        provides the first unambiguous microscopic proof of four-state
        toroidal domain control in a bulk antiferromagnet. This confirms
        that the T$^4$ topology -- the basis of FFGFT -- is not merely a
        mathematical construct but manifests as a physical domain structure
        in real materials.
  \item \textbf{The $\mathbb{Z}_2\times\mathbb{Z}_2$ structure is exact:}
        The four domains are precisely related by time-inversion and a
        spatial screw operation -- the same group that labels T$^4$ winding
        modes in FFGFT.
  \item \textbf{Purely imaginary structure factors = static T$^4$:}
        The fact that $M_\perp$ is purely imaginary (antiferromagnetic,
        head-to-tail configuration) means $\mathcal{P}''=0$: no spontaneous
        polarization is created. This is the FFGFT static ground state.
  \item \textbf{Non-volatility = topological stability:} The memory is
        non-volatile because the winding number is topologically protected.
        In FFGFT, this is the same mechanism that stabilizes the lepton mass
        hierarchy: the winding numbers $p_e=3/2$, $p_\mu=2$, $p_\tau=5/2$
        are topologically protected (Dok.~190, P2).
\end{itemize}

\noindent\textbf{Status of the interpretations in this document.}
Sections~1--5 (individual formulas) are \emph{reduction layer}: structural
analogies between the paper's formalism and FFGFT, without a new
quantitative prediction. Sections~6 (4-domain = T$^4$ modes) and~7
(non-volatility = topological stability) are \emph{bridges}: the
correspondence is algebraically exact ($\mathbb{Z}_2\times\mathbb{Z}_2$
group structure, topological charge conservation), derived from both
frameworks independently. Section~8 ($\xi$-controlled scale hierarchy)
is \emph{reduction layer}: the structure is plausible but the exact power
of $\xi$ is not determined.

% ============================================================================
\section{Fractal Path Lengthening vs.\ Tired Light -- a Necessary Distinction}
% ============================================================================

The FFGFT explanation of redshift as \emph{fractal path lengthening} in a
static universe (Doc.~190, R14) must be distinguished from the historical
\emph{Tired Light} hypothesis (Zwicky 1929), which is rejected in mainstream
astrophysics for three reasons. FFGFT avoids all three:

\begin{enumerate}
  \item \textbf{No blurring of distant sources.} Tired Light requires photon
        scattering off intervening matter, which would broaden wavefronts and
        blur distant galaxies. Fractal path lengthening involves no scattering:
        the photon follows the topological geometry of the static vacuum (the
        ``head-to-tail'' winding structure demonstrated in
        LiNi$_{0.8}$Fe$_{0.2}$PO$_4$), preserving full wavefront coherence.
  \item \textbf{Consistent geometric time dilation $(1+z)$.} Tired Light
        predicts no time dilation of supernova light curves, which is
        observationally ruled out. In FFGFT, every signal traversing a
        fractally lengthened path carries the same wavelength ratio
        $\lambda_b/\lambda_e$ in wavelength as in pulse duration -- both
        scale geometrically through the fractal path lengthening, not through
        energy loss. The ratio is preserved, and the observed supernova time
        dilation follows as a direct consequence.
        (The symbol $z = \lambda_b/\lambda_e - 1$ is a label for the
        measured ratio, not a primary FFGFT quantity, Doc.~190, R14.)
  \item \textbf{Preservation of the CMB blackbody spectrum.} Photon scattering
        would distort the spectral shape. Fractal path lengthening scales all
        frequencies uniformly ($\Delta\nu/\nu = \mathrm{const}$), preserving
        the perfect blackbody form of the CMB spectrum.
\end{enumerate}

\noindent The ferrotoroidic vacuum structure demonstrated in the paper
provides the physical mechanism for (1): the ``head-to-tail'' antiferromagnetic
configuration (purely imaginary $M_\perp$, hence $\mathcal{P}''=0$, no
spontaneous polarization) ensures coherent, scattering-free propagation.

% ============================================================================
\section{Topological Structure Across Three Scales}
% ============================================================================

\begin{table}[htbp]
\centering
\caption{The same $\mathbb{Z}_2\times\mathbb{Z}_2$ topological structure
appears at three distinct scales.}
\label{tab:scales}
\begin{tabular}{llll}
\toprule
\textbf{Scale} & \textbf{System} & \textbf{Topology} & \textbf{States} \\
\midrule
Atomic ($10^{-10}$\,m) &
  LiNi$_{0.8}$Fe$_{0.2}$PO$_4$ &
  Toroidal spin rings &
  4 domains (SNP) \\
Sub-Planck ($10^{-38}$\,m) &
  FFGFT vacuum (T$^4$) &
  Toroidal winding modes &
  4 winding classes \\
Cosmological (Mpc) &
  Observable universe &
  Fractal path lengthening &
  $H_0 = 66.2$\,km/s/Mpc (no $\Lambda$) \\
\bottomrule
\end{tabular}
\end{table}

\noindent The cosmological row uses $H_0$ from $\xi$ (Doc.~263/026) as the
geometric identification, not an equation-of-state parameter $w$: FFGFT has
no dark energy field and no expansion, so there is no $w$ to assign. The
scale invariance of the topological structure is structural, not numerical.

% ============================================================================
\section{Correspondence Table: Paper $\leftrightarrow$ FFGFT}
% ============================================================================

\begin{table}[htbp]
\centering
\caption{Verified correspondences between Qureshi et al.\ (2026) and FFGFT.
Only entries with an explicit derivation or structural argument in both
frameworks are listed. Speculative or numerically unsupported identifications
are excluded.}
\label{tab:correspondence}
\begin{tabular}{ll}
\toprule
\textbf{Nature paper (Qureshi et al.\ 2026)} & \textbf{FFGFT (T$^4$ geometry)} \\
\midrule
Head-to-tail spin rings & Closed winding loops on T$^4$ \\
Toroidal moment $\mathbf{T}\parallel\hat{\mathbf{c}}$ &
  Winding flux in compact $c$-direction \\
Spontaneous spin tilt ($\varphi$) &
  Decompactification projection ($\cos\varphi$, Layer~2) \\
4 discrete domains &
  $\mathbb{Z}_2\times\mathbb{Z}_2$ winding mode classes \\
Time inversion $1'$: $a\leftrightarrow b$ &
  Temporal winding direction ($\pm\tilde{T}$) \\
Screw axis $2_1\|y$: $'\leftrightarrow''$ &
  Spatial winding parity ($\pm c$) \\
Polarization matrix $\mathcal{P}$ (rotation) &
  Recursion $\mathcal{R}$ on T$^4$ mode vector (Doc.~203) \\
$\mathcal{P}''=0$ (purely imaginary $M_\perp$) &
  Static T$^4$ ground state (no mode creation) \\
Non-volatile memory &
  Topological charge conservation (Doc.~204) \\
Measurement ratio $P_{fi}$ (dimensionless) &
  Layer-1 ratio (no SI unit needed) \\
\bottomrule
\end{tabular}
\end{table}

\noindent\textbf{Reference.} N.~Qureshi et al., \textit{Toroidicity as a
route towards non-volatile quaternary memory in antiferromagnets},
\textit{Nature Communications} \textbf{17}, 4033 (2026).
DOI:~10.1038/s41467-026-70767-8.
